\section{Introduction}\label{sec:introduction}

A manufacturer responding to stronger demand for industrial automation must
decide how much capacity to add and which activities to expand.
Its investment changes the margins available to competing manufacturers,
whose responses in turn affect the attractiveness of that expansion.
Adjusting capacity takes time, while news of its cash-flow consequences can
change firms' valuations before the operating response is complete.
These linked decisions change firms' positions relative to their competitors
and create a connection between investment, financial valuation, and the
information reported about firms.

Firms engage in several activities at once, and their competitive positions
depend on the composition of those activities as well as their scale.
This allocation of capacity defines a firm's economic footprint.
Recent work by \citet{gawronsky_continuous_2026} represents firms through
distributions of financial-news embeddings, preserving differences among the
articles associated with each firm rather than reducing its coverage to a
single average.
Such distributions provide a way to compare firms' positions in information space.
Understanding how firms' economic footprints evolve requires an account of
the choices that change their activities.

Spatial and multi-market competition provide a foundation for that account.
Since \citet{hotelling_stability_1929}, models of product positioning have
linked firms' locations to the demand they serve.
Research on strategic substitutes and complements and on networked
competition shows how firms' choices affect rivals across markets
\citep{bulow_geanakoplos_klemperer_multimarket_1985,bimpikis_ehsani_ilkilic_cournot_2019}.
The multi-contest framework of \citet{xu_zenou_zhou_conflict_2022} formalizes
such interdependence when players allocate resources across overlapping contests.
We adapt that framework to firms allocating operating capacity across
activities that serve related customer demands, where each firm's expansion
changes the contribution margins available to its competitors.

Firms differ in the activities they can pursue at low cost, while increasing
capacity brings additional expense.
We establish conditions for a unique Nash equilibrium and derive how
opportunity and cost changes affect capacity allocations, business mix, and
rivals' responses.
To study repositioning over time, we introduce costs of changing inherited capacity.
Each period's operating choices form a Nash equilibrium conditional on the
previous allocation, so adjustment unfolds gradually as competitors respond
to one another.
Investors value the cash flows generated along this path at a constant discount factor.
They incorporate anticipated changes in cash flows into firms'
valuations immediately, even as operating capacity adjusts over subsequent
periods.

In a restricted two-opportunity economy, firms with greater initial exposure
experience larger return responses to positive exogenous free cash flow shocks.
The shock improves one opportunity relative to the other, inducing firms to
reallocate capacity toward it.
If news coverage reflects the relative allocation of operating capacity, a
firm increasing its exposure becomes closer to the initial news distributions
of more-exposed competitors and farther from those of less-exposed competitors.
With fixed reference distributions and the stated conditions on news timing
and observation errors, these responses imply a nonnegative expected local
association between relative financial returns and decreases in the target
firm's distance to other firms.
This result gives competitive repositioning an observable implication in
information space.

Financial news is itself produced through economic choices.
Investors seek information about material changes in investment, products,
and cash-flow prospects, while publishers and corporate disclosures influence
which developments receive coverage \citep{gentzkow_media_bias_2006}.
Following research that uses corporate text to study product-market
relationships \citep{hoberg_phillips_2016}, we represent a firm's articles as
a mixture of distributions associated with its activities.
The relative weight of each activity depends on both operating capacity and
the intensity of coverage.
When coverage per unit of capacity is equal across activities, news
composition follows the firm's business mix.
Changes in coverage can also alter the reported distribution while capacity
remains unchanged.

We measure differences between firms' article distributions using squared
Wasserstein distance.
Optimal transport finds a minimum-cost matching between the distributions of
article embeddings, capturing changes in their composition even when their
average vectors change little.
For each firm, we measure whether its annual news distribution moves closer
to or farther from other firms' distributions fixed at their 2018 positions.
We compare these changes with the reference firms' daily excess returns,
summed with weights determined by the proportion of the target firm's
articles published after each return date.
Centering both measures across reference firms within a target-year produces
an estimator of the predicted relative association.

Multi-agent simulations assess this estimator using operating paths and
return responses generated by the model.
Under equal reporting intensity, the population coefficient is positive in
each of sixteen simulated news histories, yet only
\ppnum[1]{ir-validation-common-1-correct-pct}\% of article samples at
observed coverage yield a numerically resolved positive estimate.
Increasing article coverage improves recovery, although measurement error
remains substantial.
In a separate observation counterfactual, changing the relative coverage of
activities also produces positive alignment while operating capacities are held fixed.
These experiments show how article sampling and the production of news affect
the ability to observe competitive adjustment.

We apply the estimator to financial news and returns for \ppnum[0]{ir-firms}
Nasdaq-listed firms from 2018 to 2022.
Across four specifications, firms' movements relative to fixed reference
distributions are positively associated with those reference firms' relative
weighted excess returns, although the relationship is small.
The evidence is descriptive, while the simulations establish why changes in
news cannot be interpreted as changes in operating capacity alone.
Together, the theory and empirical analysis extend the study of
distribution-valued firms from comparisons of their reported economic
footprints to an economic account of why those footprints may change over time.

The remainder of the paper proceeds as follows.
Section~\ref{sec:literature} reviews the related literature.
Section~\ref{sec:model} develops the economic model, its directional
prediction, the estimator and the restricted-form hypothesis.
Section~\ref{sec:application} evaluates the estimator and its measurement
uncertainty through multi-agent simulations, then describes the corporate data
and empirical design.
Section~\ref{sec:results} presents the results and their economic interpretation.
Section~\ref{sec:conclusion} concludes.
